\pagestyle{empty}
\tight
\begin{tblr}{width=\textwidth,colspec={|Q[c,m,1.9cm]|X[l,m]|},hlines,vlines,hspan=minimal,row{1}={bg=headblue},rowsep=1pt,colsep=3pt}
\SetCell{c,m}\hfil\logo\hfil & \textbf{POLITEKNIK NEGERI MALANG}\par\textbf{JURUSAN TEKNOLOGI INFORMASI}\par\textbf{PROGRAM STUDI: D4 TEKNIK INFORMATIKA} \\
\SetCell[c=2]{c} \textbf{RENCANA PEMBELAJARAN SEMESTER (RPS)} & \\
\end{tblr}
\begin{longtblr}{width=\textwidth,colspec={|Q[l,m,1.9cm]|X[0.85,l,m]|X[1.45,l,m]|X[0.75,l,m]|X[1.15,l,m]|},hlines,vlines,hspan=minimal,rowsep=1pt,colsep=2pt,row{1,3}={bg=headgray,font=\bfseries}}
MATA KULIAH & KODE & BOBOT (SKS)/jam & SEMESTER & TGL. PENYUSUNAN \\
Pemrograman Berbasis Objek & RTI253007 & 2 SKS/4 jam & 3 & 1 Juli 2025 \\
OTORISASI & Dosen Pengembang RPS & Koordinator RMK & \SetCell[c=2]{l} Ka PRODI &  \\
 & Agung Nugroho Pramudhita, S.T., M.T. &  & \SetCell[c=2]{l}{\vspace{8mm}\par Dr. Ely Setyo Astuti, S.T., M.T.} &  \\
\SetCell[r=12]{l} Capaian Pembelajaran (CP) & \SetCell[c=4]{l,bg=headgray,font=\bfseries} Capaian Pembelajaran Lulusan Program Studi (CPL-Prodi) &  &  &  \\
 & CPL02 & \SetCell[c=3]{l} Mampu menerapkan prinsip rekayasa sistem dan perangkat lunak untuk menghasilkan solusi aplikasi multi-platform sesuai kebutuhan. &  &  \\
 & CPL07 & \SetCell[c=3]{l} Mampu mengembangkan sistem komputasi cerdas dan melakukan analisis data berbasis AI untuk pengambilan keputusan. &  &  \\
 & CPL09 & \SetCell[c=3]{l} Mampu memahami cara kerja sistem komputer dan menerapkan pendekatan komputasi untuk problem solving. &  &  \\
 & \SetCell[c=4]{l,bg=headgray,font=\bfseries} Capaian Pembelajaran mata kuliah (CPL-MK) &  &  &  \\
 & CPMK0210 & \SetCell[c=3]{l} Mahasiswa mampu memahami dan menerapkan prinsip fundamental rekayasa perangkat lunak dan OOP untuk menghasilkan spesifikasi teknis yang terstruktur. &  &  \\
 & CPMK0704 & \SetCell[c=3]{l} Mahasiswa mampu menggunakan berbagai bahasa dan paradigma pemrograman berorientasi objek untuk menyelesaikan masalah komputasi. &  &  \\
 & CPMK0903 & \SetCell[c=3]{l} Mahasiswa mampu menerapkan logika komputasi, pemrograman berbasis objek, dan algoritma untuk membangun aplikasi perangkat lunak. &  &  \\
 & \SetCell[c=4]{l,bg=headgray,font=\bfseries} Kemampuan akhir tiap tahapan belajar (Sub-CPMK) &  &  &  \\
 & SCPMK0210-02501 & \SetCell[c=3]{l} Mahasiswa mampu menjelaskan prinsip OOP (enkapsulasi, inheritance, polymorfisme) dan menerapkannya dalam rancangan kelas. &  &  \\
 & SCPMK0704-02502 & \SetCell[c=3]{l} Mahasiswa mampu mengimplementasikan konsep OOP menggunakan Java dalam berbagai paradigma pemrograman. &  &  \\
 & SCPMK0903-02503 & \SetCell[c=3]{l} Mahasiswa mampu membangun aplikasi Java berbasis GUI dan database menggunakan prinsip OOP secara menyeluruh. &  &  \\
\SetCell[r=2]{l} Penilaian & \SetCell[c=4]{l} *Nilai CPMK didapat dari agregasi nilai SUB-CPMK yang dibebankan &  &  &  \\
 & \SetCell[c=4]{l}{\tiny\begin{tblr}{width=\linewidth,colspec={|Q[c,m,0.1\linewidth]|Q[c,m,0.16\linewidth]|X[l,m]|X[l,m]|X[l,m]|X[l,m]|X[l,m]|X[l,m]|Q[c,m,0.08\linewidth]|},hlines,vlines,hspan=minimal,rowsep=1pt,colsep=0.7pt,row{1,2}={bg=headgray,font=\bfseries}}
Kode CPMK & Kode SCPMK & \SetCell[c=6]{c} Bobot per Bentuk Penilaian & & & & & & Total Bobot \\
 & & Tugas & Kuis 1 & UTS & CM & Kuis 2 & UAS & \\
CPMK0210 & SCPMK0210-02501 & 5\%: prinsip OOP dan rancangan kelas & 5\%: konsep OOP & 10\%: prinsip OOP & 10\%: implementasi prinsip OOP & 0\% & 0\% & 30\% \\
CPMK0704 & SCPMK0704-02502 & 5\%: implementasi Java OOP & 0\% & 0\% & 30\%: implementasi Java berbagai paradigma & 5\%: OOP lanjut & 0\% & 40\% \\
CPMK0903 & SCPMK0903-02503 & 5\%: aplikasi GUI dan database & 0\% & 0\% & 15\%: aplikasi Java terpadu & 0\% & 10\%: UAS aplikasi Java & 30\% \\
\SetCell[c=8]{r}\textbf{Total Per Penilaian} & & & & & & & & \textbf{100\%} \\
\end{tblr}} &  &  &  \\
Diskripsi Singkat Mata Kuliah & \SetCell[c=4]{l} Mata kuliah Pemrograman Berbasis Objek membangun kemampuan merancang dan mengimplementasikan program berorientasi objek menggunakan Java, mencakup enkapsulasi, inheritance, polymorfisme, GUI, dan koneksi database. &  &  &  \\
Bahan Kajian / Materi Pembelajaran & \SetCell[c=4]{l}\begin{enumerate}\item Konsep dasar OOP: kelas, objek, enkapsulasi\item Relasi kelas, inheritance, dan overriding\item Interface, kelas abstrak, dan polymorfisme\item GUI dan Java API\item Java GUI dengan database\end{enumerate} &  &  &  \\
\SetCell[r=2]{l} Pustaka & \SetCell[c=4]{l} \textbf{Utama:}\begin{enumerate}\item Deitel, P. \& Deitel, H. 2017. Java How to Program. Pearson.\item Bloch, J. 2018. Effective Java. Addison-Wesley.\item Eckel, B. 2006. Thinking in Java. Prentice Hall.\end{enumerate} &  &  &  \\
 & \SetCell[c=4]{l} \textbf{Pendukung:} Materi LMS, dokumentasi resmi Java API, dan jobsheet praktikum. &  &  &  \\
Media Pembelajaran & \SetCell[c=2]{l} \textbf{Software:} JDK, IDE (Eclipse/IntelliJ), NetBeans, LMS. &  & \SetCell[c=2]{l} \textbf{Hardware:} Komputer/laptop, proyektor. &  \\
Nama Dosen Pengampu & \SetCell[c=4]{l} Tim Pengajar Program Studi D4 TI &  &  &  \\
Matakuliah Syarat & \SetCell[c=4]{l} RTI251006 Dasar Pemrograman &  &  &  \\
\end{longtblr}
\begin{longtblr}{width=\textwidth,colspec={|Q[c,m,0.06\textwidth]|X[1.35,l,m]|X[1.08,l,m]|X[1.16,l,m]|X[0.7,l,m]|X[1.24,l,m]|X[1.06,l,m]|X[0.98,l,m]|Q[c,m,0.05\textwidth]|},hlines,vlines,hspan=minimal,rowhead=2,row{1,2}={bg=headgreen,font=\bfseries},rowsep=1pt,colsep=1.5pt}
Mgg.\par Ke & Kemampuan Akhir Yang Direncanakan (Sub-CP-MK) & Bahan kajian (Materi Pembelajaran) & Bentuk dan Metode Pembelajaran & Estimasi Waktu & Pengalaman Belajar Mahasiswa & Kriteria \& Bentuk Penilaian & Indikator Penilaian & Bobot Penilaian (\%) \\
(1) & (2) & (3) & (4) & (5) & (6) & (7) & (8) & (9) \\
1 & SCPMK0210-02501 & Pengantar konsep PBO. & Modalitas: Blended Learning. Bentuk: Luring/Daring. Metode: Case Method, praktik pemrograman, studi kasus, dan peer review. & 1 x 2 x 50' tatap muka; 1 x 2 x 50' tugas/praktik mandiri. & Mahasiswa mengerjakan aktivitas terarah untuk pengantar konsep PBO dan menunjukkan evidence hasil belajar. & Bentuk penilaian: Pengantar Konsep PBO. Mengacu rubrik RTM. & Ketepatan konsep; kualitas implementasi; validasi dan dokumentasi. & 3 \\
2 & SCPMK0210-02501 & Kelas dan obyek. & Modalitas: Blended Learning. Bentuk: Luring/Daring. Metode: Case Method, praktik pemrograman, studi kasus, dan peer review. & 1 x 2 x 50' tatap muka; 1 x 2 x 50' tugas/praktik mandiri. & Mahasiswa mengerjakan aktivitas terarah untuk kelas dan obyek dan menunjukkan evidence hasil belajar. & Bentuk penilaian: Kelas dan Obyek. Mengacu rubrik RTM. & Ketepatan konsep; kualitas implementasi; validasi dan dokumentasi. & 3 \\
3 & SCPMK0210-02501 & Enkapsulasi. & Modalitas: Blended Learning. Bentuk: Luring/Daring. Metode: Case Method, praktik pemrograman, studi kasus, dan peer review. & 1 x 2 x 50' tatap muka; 1 x 2 x 50' tugas/praktik mandiri. & Mahasiswa mengerjakan aktivitas terarah untuk enkapsulasi dan menunjukkan evidence hasil belajar. & Bentuk penilaian: Enkapsulasi. Mengacu rubrik RTM. & Ketepatan konsep; kualitas implementasi; validasi dan dokumentasi. & 3 \\
4 & SCPMK0210-02501 & Relasi kelas. & Modalitas: Blended Learning. Bentuk: Luring/Daring. Metode: Case Method, praktik pemrograman, studi kasus, dan peer review. & 1 x 2 x 50' tatap muka; 1 x 2 x 50' tugas/praktik mandiri. & Mahasiswa mengerjakan aktivitas terarah untuk relasi kelas dan menunjukkan evidence hasil belajar. & Bentuk penilaian: Relasi Kelas. Mengacu rubrik RTM. & Ketepatan konsep; kualitas implementasi; validasi dan dokumentasi. & 3 \\
5 & SCPMK0210-02501 & Kuis 1. & Modalitas: Blended Learning. Bentuk: Luring/Daring. Metode: Case Method, praktik pemrograman, studi kasus, dan peer review. & 1 x 2 x 50' tatap muka; 1 x 2 x 50' tugas/praktik mandiri. & Mahasiswa mengerjakan aktivitas terarah untuk kuis 1 dan menunjukkan evidence hasil belajar. & Bentuk penilaian: Kuis 1. Mengacu rubrik RTM. & Ketepatan konsep; kualitas implementasi; validasi dan dokumentasi. & 5 \\
6 & SCPMK0210-02501 & Inheritance 1. & Modalitas: Blended Learning. Bentuk: Luring/Daring. Metode: Case Method, praktik pemrograman, studi kasus, dan peer review. & 1 x 2 x 50' tatap muka; 1 x 2 x 50' tugas/praktik mandiri. & Mahasiswa mengerjakan aktivitas terarah untuk inheritance 1 dan menunjukkan evidence hasil belajar. & Bentuk penilaian: Inheritance 1. Mengacu rubrik RTM. & Ketepatan konsep; kualitas implementasi; validasi dan dokumentasi. & 3 \\
7 & SCPMK0210-02501 & Inheritance 2. & Modalitas: Blended Learning. Bentuk: Luring/Daring. Metode: Case Method, praktik pemrograman, studi kasus, dan peer review. & 1 x 2 x 50' tatap muka; 1 x 2 x 50' tugas/praktik mandiri. & Mahasiswa mengerjakan aktivitas terarah untuk inheritance 2 dan menunjukkan evidence hasil belajar. & Bentuk penilaian: Inheritance 2. Mengacu rubrik RTM. & Ketepatan konsep; kualitas implementasi; validasi dan dokumentasi. & 3 \\
8 & SCPMK0210-02501 & UTS. & Modalitas: Blended Learning. Bentuk: Luring/Daring. Metode: Case Method, praktik pemrograman, studi kasus, dan peer review. & 1 x 2 x 50' tatap muka; 1 x 2 x 50' tugas/praktik mandiri. & Mahasiswa mengerjakan aktivitas terarah untuk UTS dan menunjukkan evidence hasil belajar. & Bentuk penilaian: UTS. Mengacu rubrik RTM. & Ketepatan konsep; kualitas implementasi; validasi dan dokumentasi. & 10 \\
9 & SCPMK0704-02502 & Konsep overriding dan overloading. & Modalitas: Blended Learning. Bentuk: Luring/Daring. Metode: Case Method, praktik pemrograman, studi kasus, dan peer review. & 1 x 2 x 50' tatap muka; 1 x 2 x 50' tugas/praktik mandiri. & Mahasiswa mengerjakan aktivitas terarah untuk konsep overriding dan overloading dan menunjukkan evidence hasil belajar. & Bentuk penilaian: Konsep Overriding dan Overloading. Mengacu rubrik RTM. & Ketepatan konsep; kualitas implementasi; validasi dan dokumentasi. & 7 \\
10 & SCPMK0704-02502 & Kelas abstrak. & Modalitas: Blended Learning. Bentuk: Luring/Daring. Metode: Case Method, praktik pemrograman, studi kasus, dan peer review. & 1 x 2 x 50' tatap muka; 1 x 2 x 50' tugas/praktik mandiri. & Mahasiswa mengerjakan aktivitas terarah untuk kelas abstrak dan menunjukkan evidence hasil belajar. & Bentuk penilaian: Kelas Abstrak. Mengacu rubrik RTM. & Ketepatan konsep; kualitas implementasi; validasi dan dokumentasi. & 7 \\
11 & SCPMK0704-02502 & Interface. & Modalitas: Blended Learning. Bentuk: Luring/Daring. Metode: Case Method, praktik pemrograman, studi kasus, dan peer review. & 1 x 2 x 50' tatap muka; 1 x 2 x 50' tugas/praktik mandiri. & Mahasiswa mengerjakan aktivitas terarah untuk interface dan menunjukkan evidence hasil belajar. & Bentuk penilaian: Interface. Mengacu rubrik RTM. & Ketepatan konsep; kualitas implementasi; validasi dan dokumentasi. & 7 \\
12 & SCPMK0704-02502 & Polymorfisme. & Modalitas: Blended Learning. Bentuk: Luring/Daring. Metode: Case Method, praktik pemrograman, studi kasus, dan peer review. & 1 x 2 x 50' tatap muka; 1 x 2 x 50' tugas/praktik mandiri. & Mahasiswa mengerjakan aktivitas terarah untuk polymorfisme dan menunjukkan evidence hasil belajar. & Bentuk penilaian: Polymorfisme. Mengacu rubrik RTM. & Ketepatan konsep; kualitas implementasi; validasi dan dokumentasi. & 7 \\
13 & SCPMK0704-02502 & Kuis 2. & Modalitas: Blended Learning. Bentuk: Luring/Daring. Metode: Case Method, praktik pemrograman, studi kasus, dan peer review. & 1 x 2 x 50' tatap muka; 1 x 2 x 50' tugas/praktik mandiri. & Mahasiswa mengerjakan aktivitas terarah untuk kuis 2 dan menunjukkan evidence hasil belajar. & Bentuk penilaian: Kuis 2. Mengacu rubrik RTM. & Ketepatan konsep; kualitas implementasi; validasi dan dokumentasi. & 5 \\
14 & SCPMK0903-02503 & GUI. & Modalitas: Blended Learning. Bentuk: Luring/Daring. Metode: Case Method, praktik pemrograman, studi kasus, dan peer review. & 1 x 2 x 50' tatap muka; 1 x 2 x 50' tugas/praktik mandiri. & Mahasiswa mengerjakan aktivitas terarah untuk GUI dan menunjukkan evidence hasil belajar. & Bentuk penilaian: GUI. Mengacu rubrik RTM. & Ketepatan konsep; kualitas implementasi; validasi dan dokumentasi. & 10 \\
15 & SCPMK0903-02503 & Java API. & Modalitas: Blended Learning. Bentuk: Luring/Daring. Metode: Case Method, praktik pemrograman, studi kasus, dan peer review. & 1 x 2 x 50' tatap muka; 1 x 2 x 50' tugas/praktik mandiri. & Mahasiswa mengerjakan aktivitas terarah untuk Java API dan menunjukkan evidence hasil belajar. & Bentuk penilaian: Java API. Mengacu rubrik RTM. & Ketepatan konsep; kualitas implementasi; validasi dan dokumentasi. & 10 \\
16 & SCPMK0903-02503 & Java GUI dan database. & Modalitas: Blended Learning. Bentuk: Luring/Daring. Metode: Case Method, praktik pemrograman, studi kasus, dan peer review. & 1 x 2 x 50' tatap muka; 1 x 2 x 50' tugas/praktik mandiri. & Mahasiswa mengerjakan aktivitas terarah untuk Java GUI dan database dan menunjukkan evidence hasil belajar. & Bentuk penilaian: Java GUI dan Database. Mengacu rubrik RTM. & Ketepatan konsep; kualitas implementasi; validasi dan dokumentasi. & 10 \\
17 & SCPMK0704-02502, SCPMK0903-02503 & UAS. & Modalitas: Blended Learning. Bentuk: Luring/Daring. Metode: Case Method, praktik pemrograman, studi kasus, dan peer review. & 1 x 2 x 50' tatap muka; 1 x 2 x 50' tugas/praktik mandiri. & Mahasiswa mengerjakan aktivitas terarah untuk UAS dan menunjukkan evidence hasil belajar. & Bentuk penilaian: UAS. Mengacu rubrik RTM. & Ketepatan konsep; kualitas implementasi; validasi dan dokumentasi. & 4 \\
\end{longtblr}
